% 05_theory.tex  --  Theorem A WRITTEN; B, C structured stubs
\section{Theoretical Analysis}
\label{sec:theory}

This section analyzes LA-JSSA. Because a single tractable theorem cannot
simultaneously capture the learned detector, the continuous control, and the
finite-time learning dynamics, we analyze the system through a deliberate
\emph{hierarchy of abstractions}---true system (Sec.~\ref{sec:system_model})
$\to$ belief map (Sec.~\ref{sec:framework}) $\to$ fixed-policy evaluation
(Thms.~\ref{thm:tradeoff},~\ref{thm:cvar}) $\to$ tabular surrogate
(Prop.~\ref{thm:regret})---each tier exposing one structural property while
holding the others fixed. The results answer three \emph{distinct} questions
about the same system, and we state up front what each does and does
\emph{not} claim:
\begin{itemize}
\item \textbf{Operating-point analysis} (Theorems~\ref{thm:tradeoff},
\ref{thm:cvar}): a \emph{fixed-policy evaluation} of how the belief threshold
$\thr$---for a fixed action geometry---trades sensing error against
throughput and PU protection, under the mean and the CVaR constraint. These
characterize a static operating point; they are not statements about the
learning process.
\item \textbf{Representation analysis} (Proposition~\ref{thm:sufficiency}):
an information-theoretic comparison of the belief-map state against a coarser
fixed-channel state \emph{under a common action space}. It concerns what is
\emph{representable}, not how it is learned.
\item \textbf{Learning analysis} (Proposition~\ref{thm:regret}): a finite-time
regret bound for a \emph{tabular surrogate} of the access problem. It gives
qualitative insight into the structure of the learning cost; it is
\emph{not} a guarantee for the deep agent of Section~\ref{sec:framework},
which serves as a heuristic solver of the surrogate.
\end{itemize}
We make no claim that the deep RL algorithm inherits these guarantees; the
theory characterizes the surrogate hierarchy, and the experiments
(Section~\ref{sec:simulation}) test whether the predicted \emph{qualitative}
structure carries over to the deep agent.

\begin{remark}[Relation to classical sensing--throughput trade-off]
The classical sensing--throughput trade-off \cite{liang2008sensing} optimizes
the \emph{sensing duration} and an energy threshold for a \emph{single}
channel along the time axis. Theorem~\ref{thm:tradeoff} differs in kind: the
lever is the belief threshold $\thr$ on a \emph{two-dimensional} continuous
TF occupancy estimate, coupling the localization error of a learned detector
to the \emph{access} reward through the collision event
\eqref{eq:collision}.
\end{remark}

\subsection{Assumptions}
\label{subsec:assumptions}

The access stage consumes the output of a \emph{time--frequency localizer},
not a per-channel energy detector. A key structural feature of the
sense-then-access frame simplifies what matters: within one frame, the
sensing phase estimates which \emph{frequency} ranges are unoccupied, and the
access phase of the \emph{same} frame reuses those ranges. The temporal
support of an emission is thus fixed by the frame boundary and given to both
phases; it is the \emph{frequency} extent of each emission that the localizer
must pin down, since a frequency-localization error is what causes the access
phase to either collide (under-estimated occupancy) or waste spectrum
(over-estimated occupancy). We therefore measure localization accuracy by the
\emph{frequency-interval} overlap and relate it to the cell-level error that
drives the access analysis, so that $\Pmd,\Pfa$ are \emph{derived} from
localization quality rather than postulated.

\begin{assumption}[Frequency-localization error model]
\label{as:localization}
For each true PU emission, let its occupied frequency interval in a frame be
$\fint_k=[\flo^{k},\fhi^{k}]$, and let the localizer emit the estimate
$\hat\fint_k=[\hat\flo^{k},\hat\fhi^{k}]$. Measure the frequency-localization
accuracy by the one-dimensional overlap
\begin{equation}
\iou_k\defeq\frac{|\fint_k\cap\hat\fint_k|}{|\fint_k\cup\hat\fint_k|}\in[0,1],
\qquad
\eloc\defeq 1-\E[\iou_k],
\label{eq:freq_iou}
\end{equation}
the mean frequency-localization error. After rasterizing $\hat\fint_k$ to the
$\Nf$ frequency cells and thresholding the belief at $\thr$, the induced
per-cell miss and false-alarm rates satisfy
\begin{equation}
\Pmd(\thr)\le \cmd\,\eloc + \rmd(\thr),
\qquad
\Pfa(\thr)\le \cfa\,\eloc + \rfa(\thr),
\label{eq:loc_to_cell}
\end{equation}
where $\rmd,\rfa$ are the residual detection errors of a perfectly localized
interval---the contribution of the \emph{confidence threshold} $\thr$ alone,
i.e.,\ where the soft score is cut---and $\cmd,\cfa>0$ are geometry
constants scaling the \emph{localization} contribution. Thus tightening the
frequency localization (higher IoU, smaller $\eloc$) lowers both cell-level
error floors, and adjusting $\thr$ trades $\Pmd$ against $\Pfa$ along the
residual receiver operating characteristic (ROC). Temporal-support error does
not appear: it is absorbed by the
frame alignment shared between the sensing and access phases.
\end{assumption}

\begin{assumption}[Cell-level error model]
\label{as:cellerror}
Conditioned on the true occupancy, the thresholded belief decisions
$\hat{\occ}_{m,n}\defeq\Ind\{\bel_{m,n}\ge\thr\}$ across cells are
characterized by a common miss-detection probability $\Pmd(\thr)$ and
false-alarm probability $\Pfa(\thr)$ as in \eqref{eq:pmd}--\eqref{eq:pfa},
each bounded via \eqref{eq:loc_to_cell} of
Assumption~\ref{as:localization}. Both are continuous and monotonic in
$\thr$: $\Pmd$ is non-decreasing and $\Pfa$ is non-increasing in $\thr$.
\end{assumption}

\begin{assumption}[Receiver operating characteristic]
\label{as:roc}
The detector admits a continuously differentiable ROC, i.e.,\ there is a
differentiable, strictly decreasing function $\rho$ with
$\Pmd = \rho(\Pfa)$ (equivalently the detection probability
$P_{\mathrm{d}}=1-\Pmd$ is a concave, increasing function of $\Pfa$). This
holds for energy- and learning-based detectors operating above the
noise floor.
\end{assumption}

\begin{assumption}[Fixed action geometry; bounded rate and penalty]
\label{as:rate}
The analysis is conducted for a \emph{fixed} SU action geometry: the chosen
block covers $A\defeq|\mathcal{A}(\act)|$ cells, of which $A_{\mathrm{idle}}$
are truly idle and $|\mathcal{O}(\act)|$ truly occupied, and these counts are
treated as constants independent of $\thr$ (the threshold governs detection,
not block placement). A truly idle accessed cell yields normalized rate
$\bar r>0$; a transmission overlapping any occupied cell is lost and incurs a
retransmission penalty captured by $\kappa>0$.
\end{assumption}

\begin{assumption}[Single-peaked throughput]
\label{as:concave}
The throughput objective $\thru(\thr)=\bar r A_{\mathrm{idle}}(1-\Pfa(\thr))
-\kappa\Pcol(\thr)$ is continuous and \emph{quasiconcave} in $\thr$ on
$[0,1]$, so that its stationary point is a global maximizer. We restrict
attention to ROC families for which this holds---e.g.,\ those whose
$\Pfa(\thr)$ is convex and $\Pcol(\thr)$ concave (or both affine) over the
operating range, giving a single-crossing first-order condition. This covers
the central, near-linear regime of typical energy- and learning-based
detector ROCs; outside it, the result of Theorem~\ref{thm:tradeoff} holds
only as a first-order/boundary characterization (stated explicitly in the
proof).
\end{assumption}

\subsection{Error Propagation and the Three-Way Trade-Off}
\label{subsec:thmA}

We begin with two lemmas that map the cell-level error to the two access
performance measures.

\begin{lemma}[Collision bound]
\label{lem:collision}
Fix an SU action $\act$ and let $\mathcal{O}(\act)$ be the set of truly
occupied cells overlapping the chosen block. Under
Assumption~\ref{as:cellerror}, the collision probability obeys the
distribution-free union bound
\begin{equation}
\Pcol(\thr)\le\!\!\sum_{(m,n)\in\mathcal{O}(\act)}\!\!
\Prob\big(\hat\occ_{m,n}{=}0\mid \occ_{m,n}{=}1\big)
=|\mathcal{O}(\act)|\,\Pmd(\thr).
\label{eq:lem_collision_union}
\end{equation}
If, in addition, the per-cell miss events are conditionally independent given
the occupancy (or positively associated), the bound tightens to
\begin{equation}
\Pcol(\thr)\;\le\;1-\bigl(1-\Pmd(\thr)\bigr)^{|\mathcal{O}(\act)|},
\label{eq:lem_collision}
\end{equation}
which agrees with \eqref{eq:lem_collision_union} to first order,
$\Pcol(\thr)\le|\mathcal{O}(\act)|\,\Pmd(\thr)+o(\Pmd)$.
\end{lemma}

\begin{IEEEproof}
A collision occurs iff at least one cell in $\mathcal{O}(\act)$ is declared
idle. The union bound over these $|\mathcal{O}(\act)|$ events, each of
marginal probability $\Pmd(\thr)$ by Assumption~\ref{as:cellerror}, gives
\eqref{eq:lem_collision_union} with no dependence assumption. Under
conditional independence the no-collision probability factorizes as
$(1-\Pmd)^{|\mathcal{O}(\act)|}$ (and the Fortuin--Kasteleyn--Ginibre (FKG)
inequality gives the same direction under positive association), yielding \eqref{eq:lem_collision};
the first-order expansion follows from $(1-x)^{k}=1-kx+o(x)$. We use the
conservative form \eqref{eq:lem_collision_union} in the sequel, so no
counterfactual quantity is required.
\end{IEEEproof}

\begin{lemma}[Throughput bound]
\label{lem:throughput}
Under the same assumptions, the expected normalized SU throughput obeys
\begin{equation}
\thru(\thr) \;\ge\; \bar{r}\,A_{\mathrm{idle}}\bigl(1-\Pfa(\thr)\bigr)
\;-\; \kappa\,\Pcol(\thr),
\label{eq:lem_throughput}
\end{equation}
where $A_{\mathrm{idle}}$ is the number of truly idle cells available to the
block and $\kappa>0$ aggregates the retransmission penalty of
Assumption~\ref{as:rate}.
\end{lemma}

\begin{IEEEproof}
A truly idle cell contributes rate $\bar{r}$ only if it is not falsely
flagged as occupied, which occurs with probability $1-\Pfa$; summing the
$A_{\mathrm{idle}}$ idle cells gives the first term by linearity of
expectation. Each collision event removes the affected rate and adds a
retransmission cost, captured by $-\kappa\Pcol$. Lower-bounding the
(non-negative) higher-order interaction terms by zero yields
\eqref{eq:lem_throughput}.
\end{IEEEproof}

\begin{theorem}[Three-way trade-off and operating point under fixed action geometry]
\label{thm:tradeoff}
Consider maximizing the expected throughput \eqref{eq:lem_throughput} subject
to the PU-protection constraint $\Pcol(\thr)\le\Pcolmax$ over the belief
threshold $\thr\in[0,1]$. Under Assumptions~\ref{as:cellerror}--\ref{as:rate}:
\begin{enumerate}
\item[(i)] \emph{(Feasibility frontier.)} There exists a threshold
$\thr_{\min}$, the smallest $\thr$ with
$|\mathcal{O}(\act)|\,\Pmd(\thr)\le\Pcolmax$, such that the constraint is
satisfied iff $\thr\ge\thr_{\min}$.
\item[(ii)] \emph{(Trade-off surface.)} On the feasible set the achievable
throughput is governed by the map
$\thr\mapsto\bigl(\Pfa(\thr),\Pmd(\thr)\bigr)$ and, via
Assumption~\ref{as:roc}, traces a one-dimensional curve on the
sensing-accuracy/throughput/PU-protection surface.
\item[(iii)] \emph{(Optimal operating point.)} The throughput-maximizing
feasible threshold is
\begin{equation}
\thr^{\star}=
\begin{cases}
\thr_{\mathrm{u}}, & \text{if } \thr_{\mathrm{u}}\ge\thr_{\min},\\[2pt]
\thr_{\min}, & \text{otherwise,}
\end{cases}
\label{eq:thmA_opt}
\end{equation}
where $\thr_{\mathrm{u}}$ is the unconstrained stationary point solving
\begin{equation}
\bar{r}\,A_{\mathrm{idle}}\,\Pfa'(\thr)
\;=\;-\,\kappa\,\Pcol'(\thr).
\label{eq:thmA_foc}
\end{equation}
That is, the optimum is interior when the marginal throughput gained by
admitting more spectrum (lower $\Pfa$) balances the marginal collision cost;
otherwise it sits on the PU-protection boundary.
\end{enumerate}
\end{theorem}

\begin{IEEEproof}
(i) By Assumption~\ref{as:cellerror}, $\Pmd(\thr)$ is non-decreasing, hence
the right side of \eqref{eq:lem_collision} is non-decreasing in $\thr$;
\emph{decreasing} $\thr$ lowers $\Pmd$ and thus $\Pcol$. Therefore the
feasible set $\{\thr:\Pcol(\thr)\le\Pcolmax\}$ is an upper interval
$[\thr_{\min},1]$ whose left endpoint is the stated $\thr_{\min}$ by
continuity (Assumption~\ref{as:cellerror}).
(ii) Eliminating $\thr$ through the ROC relation $\Pmd=\rho(\Pfa)$
(Assumption~\ref{as:roc}) parameterizes throughput and collision jointly by
$\Pfa$, giving the claimed curve.
(iii) Substitute \eqref{eq:lem_collision} into \eqref{eq:lem_throughput} to
write throughput as $\thru(\thr)=\bar r A_{\mathrm{idle}}(1-\Pfa(\thr))
-\kappa\Pcol(\thr)$. Differentiating and setting to zero gives the
first-order condition \eqref{eq:thmA_foc}. Under
Assumption~\ref{as:concave} (single-peaked $\thru$), any interior stationary
point $\thr_{\mathrm u}$ is the global maximizer on $[0,1]$; projecting onto
the feasible interval $[\thr_{\min},1]$ yields \eqref{eq:thmA_opt}. Absent
Assumption~\ref{as:concave}, \eqref{eq:thmA_foc} remains valid as a
\emph{first-order optimality condition} and \eqref{eq:thmA_opt} as the
characterization of the constrained optimum among stationary and boundary
points, without the global-optimality claim.
\end{IEEEproof}

\begin{corollary}[Price of imperfect localization]
\label{cor:price}
Let $\thru^{\circ}=\bar r A_{\mathrm{idle}}$ be the oracle throughput under
perfect sensing ($\Pmd=\Pfa=0$). Then the throughput gap satisfies
$\thru^{\circ}-\thru(\thr^{\star})\le
\bar r A_{\mathrm{idle}}\Pfa(\thr^{\star})+\kappa\Pcolmax$, i.e.,\ the loss is
controlled linearly by the false-alarm rate at the operating point and by
the admitted collision budget.
\end{corollary}

\begin{corollary}[Access cost of localization error]
\label{cor:loc}
Combining Corollary~\ref{cor:price} with the localization-error model
\eqref{eq:loc_to_cell}, the throughput gap is controlled by the mean
localization error $\eloc=1-\E[\iou]$:
\begin{equation}
\thru^{\circ}-\thru(\thr^{\star})\le
\bar r A_{\mathrm{idle}}\bigl(\cfa\,\eloc+\rfa(\thr^{\star})\bigr)
+\kappa\Pcolmax .
\label{eq:loc_gap}
\end{equation}
Hence each unit of improvement in frequency-localization IoU translates
linearly into recovered access throughput. This is the sense in which sharper
frequency localization---the access-relevant component of time--frequency
localization---directly improves downstream access, closing the loop the
title names.
\end{corollary}

\begin{remark}[Design implication]
Once $\Pcolmax$ is fixed, Corollaries~\ref{cor:price}--\ref{cor:loc} show the
remaining throughput levers are the false-alarm rate at $\thr^{\star}$ and the
frequency-localization error $\eloc$---exactly what a tighter-interval,
higher-IoU localizer buys, quantifying the value of improved frequency
localization.
\end{remark}

\subsection{A Risk-Sensitive (CVaR) Operating Point}
\label{subsec:thmD}
Theorem~\ref{thm:tradeoff} constrains the \emph{expected} collision. A
per-frame collision is, however, a near-degenerate (Bernoulli) event on which
a tail risk measure reduces to a scaled mean; the tail only becomes
meaningful once collisions are aggregated over time. We therefore define the
\emph{collision-window variable}
\begin{equation}
\Zcol^{W}(\thr)\;\defeq\;\sum_{t=1}^{W}\mathrm{Col}_t(\thr),
\label{eq:collision_window}
\end{equation}
the number of PU collisions over a window of $W$ consecutive access
opportunities. Bursty interference---rare windows with many collisions---is
exactly what a regulator penalizes and what a mean constraint misses. The
distributional formulation of Section~\ref{subsec:dist_rl} replaces the mean
constraint with a CVaR constraint on this window
variable,
\begin{equation}
\begin{aligned}
&\CVaR_{\cvarcol}\!\left[\Zcol^{W}(\thr)\right]\le\PcolmaxW,\\
&\CVaR_{\cvarcol}[\Zcol^{W}]\defeq\frac{1}{\cvarcol}\!\int_{1-\cvarcol}^{1}\!\!
\Qinv_{\Zcol^{W}}(u)\,du,
\end{aligned}
\label{eq:cvar_constraint}
\end{equation}
where $\cvarcol\in(0,1]$ is the tail level ($\cvarcol\to1$ recovers the mean
constraint) and $\PcolmaxW$ the per-window budget. We compare the threshold
$\thrcvar$ feasible under \eqref{eq:cvar_constraint} with the $\thr^{\star}$
of Theorem~\ref{thm:tradeoff}. This window-level definition is the quantity
measured in Section~\ref{sec:simulation} (Fig.~\ref{fig:cvartail}).

\begin{assumption}[Tail-monotone collision window]
\label{as:tailmono}
The collision-window variable $\Zcol^{W}$ of \eqref{eq:collision_window} is
defined over \emph{fixed-length, non-overlapping} windows of $W$ consecutive
access opportunities, with $W$ a constant independent of the horizon. For
every fixed $\thr$, $\Zcol^{W}(\thr)$ has a continuous cumulative
distribution function (CDF), and lowering
$\thr$ increases $\Zcol^{W}$ in the usual stochastic order (more missed
detections produce stochastically more collisions per window).
\end{assumption}

\begin{theorem}[Conservativeness of the CVaR operating point]
\label{thm:cvar}
Under Assumptions~\ref{as:cellerror}--\ref{as:rate} and
\ref{as:tailmono}, for any tail level $\cvarcol\in(0,1)$ the CVaR-feasible
threshold is at least as conservative as the mean-feasible one,
\begin{equation}
\thrcvar \;\ge\; \thr^{\star}.
\label{eq:cvar_conservative}
\end{equation}
Moreover $\thrcvar\to\thr^{\star}$ as the collision return becomes
deterministic (e.g.,\ as the localizer sharpens or $\cvarcol\to1$).
\end{theorem}

\begin{IEEEproof}
Since $\CVaR_{\cvarcol}[\Zcol^{W}]\ge\E[\Zcol^{W}]$ for all
$\cvarcol\in(0,1]$ with equality iff $\Zcol^{W}$ is deterministic (a standard
coherent-risk property), the CVaR constraint \eqref{eq:cvar_constraint} is
tighter than the corresponding mean constraint. By
Assumption~\ref{as:tailmono} both $\E[\Zcol^{W}(\thr)]$ and
$\CVaR_{\cvarcol}[\Zcol^{W}(\thr)]$ are non-increasing in $\thr$, so each
constraint defines an upper feasible interval; the tighter constraint has the
larger left endpoint, giving $\thrcvar\ge\thr^{\star}$, i.e.,~\eqref{eq:cvar_conservative}.
As the window variable concentrates, the
CVaR--mean gap $\to 0$ and the two intervals coincide, giving the limit.
\end{IEEEproof}

\begin{corollary}[First-order conservativeness margin]
\label{cor:cvargap}
If, in addition, $\Pcol(\cdot)$ is differentiable at $\thr^\star$ with
$\Pcol'(\thr^\star)\neq 0$ and the mean constraint is active there, then to
first order
\begin{equation}
\thrcvar-\thr^{\star}\;\approx\;
\frac{\CVaR_{\cvarcol}[\Zcol^{W}(\thr^{\star})]-\E[\Zcol^{W}(\thr^{\star})]}
{\bigl|\Pcol'(\thr^{\star})\bigr|},
\label{eq:cvar_gap}
\end{equation}
i.e.,\ the conservativeness margin scales with the collision tail spread
divided by the collision sensitivity. The associated throughput price is at
most $\bar r A_{\mathrm{idle}}\,[\Pfa(\thrcvar)-\Pfa(\thr^{\star})]$. This is
a local (first-order) estimate and is not claimed to hold globally.
\end{corollary}

\begin{IEEEproof}
A first-order expansion of the binding CVaR constraint about $\thr^{\star}$
gives $\CVaR_{\cvarcol}[\Zcol^{W}(\thr^{\star})]
+\Pcol'(\thr^{\star})(\thrcvar-\thr^{\star})\approx\Pcolmax
=\E[\Zcol^{W}(\thr^{\star})]$, where the last equality uses that the mean
constraint binds at $\thr^{\star}$. Solving for the threshold increment and
using $\Pcol'<0$ yields \eqref{eq:cvar_gap}; the throughput price follows
from Lemma~\ref{lem:throughput} evaluated at the two thresholds.
\end{IEEEproof}

\begin{remark}[Why distributional RL]
Theorem~\ref{thm:cvar} and Corollary~\ref{cor:cvargap} show the value of
learning the full return distribution: only then is the tail quantity
$\CVaR_{\cvarcol}[\Zcol^{W}]$ estimable online, enabling the conservative
operating point that a mean estimate cannot target.
\end{remark}

\subsection{Information Refinement: Sufficiency of the Belief State}
\label{subsec:thmB}
We compare two \emph{state representations} for the \emph{same} continuous
action space of Section~\ref{subsec:access_tier}: the continuous belief map
$\state$ of \eqref{eq:state}, and any fixed channelization that summarizes
the TF plane into $K$ blocks. The comparison is purely one of information,
with the action set held fixed.

\begin{proposition}[Information refinement]
\label{thm:sufficiency}
Suppose the fixed-channel state is a deterministic (measurable) function of
the belief map, $\state_{\mathrm{dis}}=g(\state)$. Let $\PIbel$ and $\PIdis$
be the policy classes measurable w.r.t.\ $\state$ and $\state_{\mathrm{dis}}$
respectively, over the common action space. Then
$\PIdis\subseteq\PIbel$ and
\begin{equation}
\sup_{\pol\in\PIbel}\Vfun^{\pol}\;\ge\;\sup_{\pol\in\PIdis}\Vfun^{\pol}.
\end{equation}
\end{proposition}

\begin{IEEEproof}
Since $\state_{\mathrm{dis}}=g(\state)$, the $\sigma$-algebra generated by
$\state_{\mathrm{dis}}$ is contained in that generated by $\state$; any
$\state_{\mathrm{dis}}$-measurable policy is therefore $\state$-measurable,
giving $\PIdis\subseteq\PIbel$. The supremum of the value over a larger
policy class is no smaller, which is the data-processing (information
monotonicity) inequality for the control value.
\end{IEEEproof}

The premise $\state_{\mathrm{dis}}=g(\state)$ holds for the standard
constructions in which the fixed-channel occupancy is obtained by averaging
or thresholding the belief map within each grid block. We emphasize that
Proposition~\ref{thm:sufficiency} is an information-refinement statement,
not a claim that \emph{every} conceivable fixed-channel scheme is dominated;
it applies precisely to those that coarsen the belief. The next example shows
the inequality can be strict.

\begin{example}[Strict gap under grid misalignment]
\label{ex:strictgap}
Consider a single PU emission whose support straddles the interior of a grid
block under a fixed channelization, occupying a fraction $\delta\in(0,1)$ of
that block. Any policy measurable w.r.t.\ the block-level state must treat
the block as a unit: accessing it collides with probability tied to
$\delta$, while declining it forgoes the idle fraction $1-\delta$. A
belief-map policy resolves the sub-block boundary and accesses only the idle
portion. The resulting value gap is at least
$\bar r\,(1-\delta)\,A_{\mathrm{cell}}-\kappa\,\delta$, which is strictly
positive for an interval of $\delta$, establishing
$\sup_{\PIbel}\Vfun>\sup_{\PIdis}\Vfun$ on this instance.
\end{example}
\noindent\emph{Derivation in Appendix~\ref{sec:appendix_suff}.}

\subsection{A Stylized Regret Analysis Under Piecewise-Stationary Dynamics}
\label{subsec:thmC}

As stated in the scope bullets above, the following result is for a
\emph{stylized tabular surrogate}---a tile-coded, finite state--action space
under a sliding-window upper-confidence index---and exposes the qualitative
structure we expect to carry over: a learning cost vanishing in rate plus a
non-vanishing floor set by sensing error. Two regularity conditions are needed.

\begin{assumption}[Bounded rewards and finite effective decision space]
\label{as:bounded}
The per-step reward \eqref{eq:reward} is bounded in $[0,1]$ after
normalization, and the access policy acts on a finite effective
state--action space of size $\mathsf{SA}$ obtained by tile-coding the belief
state and using the discrete resource-block action set of
Section~\ref{subsec:access_tier}. Reward observations are conditionally
$\sigma$-sub-Gaussian given the state and action.
\end{assumption}

\begin{assumption}[Piecewise-stationary PU with bounded switches]
\label{as:switch}
PU occupancy follows a Markov process that is stationary on each of
$\Upsilon_{\bigT}+1$ segments separated by $\Upsilon_{\bigT}$ change points,
with every segment of length at least $L_{\min}$. Within a segment the
induced belief-state MDP is fixed.
\end{assumption}

\begin{assumption}[Stationary, policy-independent sensing error]
\label{as:stationary_sensing}
Within a segment, the detector operates at a fixed threshold $\thr^\star$ and
its per-cell error rates $(\Pmd(\thr^\star),\Pfa(\thr^\star))$ are constant
and \emph{independent of the access policy}: the localizer is a fixed sensing
channel that the agent does not adapt. Consequently the per-step value gap
between the segment-optimal belief-state policy and the perfect-sensing
oracle is stationary within the segment.
\end{assumption}

\begin{proposition}[Stylized regret: a learning term plus a policy-independent sensing gap]
\label{thm:regret}
Under Assumptions~\ref{as:cellerror}--\ref{as:rate} and
\ref{as:bounded}--\ref{as:stationary_sensing}, an access policy that runs a
sliding-window upper-confidence index over the belief state, with window
tuned to $\Upsilon_{\bigT}$, incurs regret relative to the per-segment
optimal belief-state policy bounded by
\begin{equation}
\regret(\bigT)\le
\underbrace{C_1\sqrt{\mathsf{SA}\,\bigT\,\Upsilon_{\bigT}\log\bigT}}
_{\text{learning (sublinear)}}
+
\underbrace{C_2\,\bigT\big(\Pmd(\thr^{\star}){+}\Pfa(\thr^{\star})\big)}
_{\text{sensing gap (linear)}},
\label{eq:thmC}
\end{equation}
for constants $C_1,C_2>0$ depending only on $\sigma$, $\bar r$, and $\kappa$.
The first term is $o(\bigT)$ whenever $\Upsilon_{\bigT}=o(\bigT/\log\bigT)$.
The second is a \emph{policy-independent} gap set by the fixed sensing channel
at $\thr^\star$ (Assumption~\ref{as:stationary_sensing}): no policy in the
class can reduce it, since under this model the sensing error enters the
reward as a stationary additive term rather than through the transition
kernel. We do not claim the gap is irreducible if the agent were allowed to
co-adapt the localizer, which is outside the present model.
\end{proposition}
\noindent\emph{Proof in Appendix~\ref{sec:appendix_regret}.}

\subsection{Discussion}
\label{subsec:theory_discussion}
Theorem~\ref{thm:tradeoff} converts the qualitative ``better sensing helps
access'' intuition into an operating point \eqref{eq:thmA_opt} that the
simulation of Section~\ref{sec:simulation} will reproduce, validating the
theory. Proposition~\ref{thm:sufficiency} explains \emph{why} the continuous
belief state is the right representation, and Proposition~\ref{thm:regret}
separates the learning cost from the policy-independent sensing-induced gap,
clarifying which term improved localization can shrink.
